% TOP_PROBLEM: {"problemId": "problem.martin-s-conjecture-for-definable-turing-invariant-functions", "canonicalTitle": "Martin's Conjecture for Definable Turing-Invariant Functions", "releaseRank": 415, "primaryDomain": "logic_foundations_set_theory", "primaryDomainLabel": "Logic, foundations & set theory", "displayStatus": "open", "releaseStatus": "open"}
% !TeX program = lualatex
% ATTEMPT_STATUS: unresolved
% Research continuation by GPT_6_Astra_Ultra, 27 September 2026.
\documentclass[11pt]{article}
\usepackage[margin=25mm]{geometry}
\usepackage{fontspec}
\setmainfont{Latin Modern Roman}
\usepackage{amsmath,amssymb,mathtools}
\usepackage{unicode-math}
\setmathfont{Latin Modern Math}
\usepackage{xurl,hyperref}
\hypersetup{colorlinks=true,urlcolor=blue}
\setcounter{secnumdepth}{0}
\setlength{\parindent}{0pt}
\setlength{\parskip}{5pt}
\setlength{\emergencystretch}{3em}
\begin{document}
\section*{\texorpdfstring{Martin's Conjecture for Definable Turing-Invariant Functions}{Martin's Conjecture for Definable Turing-Invariant Functions}}\label{415}
\subsection{Definitions and mathematical statement}
Assume the Axiom of Determinacy in the source's ZF setting. For subsets of ${\mathbb{N}}$, $X\le_TY$ means $X$ is computable with oracle $Y$, and $X'$ is the oracle halting problem. A cone is $\{X:X\ge_TA\}$ for some $A$. A function $f:\mathcal P({\mathbb{N}})\to\mathcal P({\mathbb{N}})$ is Turing-invariant if $X\equiv_TY$ implies $f(X)\equiv_Tf(Y)$. Put $f\le_Mg$ when $f(X)\le_Tg(X)$ throughout some cone. Martin's two assertions are: every such $f$ is constant in Turing degree on a cone or satisfies $X\le_Tf(X)$ on a cone; and the functions in the latter class, modulo cone equivalence, are well-ordered by $\le_M$, with successor $[f]\mapsto[X\mapsto f(X)']$. ``Increasing'' here means above the identity, not necessarily order-preserving. The jump-iterate slogan is expressed by this precise successor law, without inventing a convention for every transfinite iterate.
\par\smallskip\textit{Formulation and concept references:} \href{https://arxiv.org/html/2510.19147}{[S1]}.

\subsection{Short English statement}
Under determinacy, are Turing-invariant functions eventually constant or above the identity, with the nonconstant cone-equivalence classes well-ordered and successive classes obtained by the Turing jump?
\subsection{Research attempt}
\textit{GPT-6 Astra Ultra; 27 September 2026. Status: UNRESOLVED.}

I work in the stated $\mathsf{ZF}+\mathsf{AD}$ setting. The main
question is not whether familiar examples form a jump hierarchy, but
whether every Turing-invariant function eventually belongs to the
asserted structure. This attempt derives the cone decision principle,
settles functions with countable degree range, and separates strict
jump growth from the much stronger successor assertion.

\paragraph{How determinacy supplies a cone.}
Let $A\subseteq2^{\mathbb N}$ be invariant under Turing equivalence.
Play the ordinary infinite binary game in which the two players
alternately contribute bits to a real $Z$. Player I wins when $Z\in A$.
By determinacy one player has a winning strategy $\sigma$, coded by
a real. Given any real $X\ge_T\sigma$, let the other player put the
successive bits of $X$ on that player's turns, while the winning player
follows $\sigma$. The resulting play $Z$ computes $X$ from those
turns. Conversely $X$ computes the whole play, because it computes
$\sigma$ and can simulate each finite stage. Thus $Z\equiv_T X$.
If I has the winning strategy then $X\in A$ for every $X\ge_T\sigma$;
if II has it then $X\notin A$ throughout that cone. This proves the
cone theorem in the form needed here.

Two cones with bases $A,B$ contain the cone with base $A\oplus B$.
It follows that the cone-containing invariant sets form an ultrafilter:
one of a set and its complement contains a cone, never both. This is
a statement about invariant sets. An arbitrary set of reals, such as
the reals with first bit one, need not contain or avoid a cone.

\paragraph{Countable degree range is eventually constant.}
For this paragraph use countable choice for sets of reals, a standard
consequence of $\mathsf{AD}$ recorded in [E1, Section 1.4]. No axiom
of dependent choice beyond the stated setting is inserted. If
$C_j$ is a sequence of cones, choose bases $A_j$ and take their
countable join $A=\bigoplus_j A_j$. Then every real above $A$ is in
every $C_j$. Hence the cone ultrafilter is countably complete.

Suppose that on a cone the values of $f$ have degrees among a given
sequence $\deg_T(B_0),\deg_T(B_1),\ldots$. The sets
\[
 E_j=\{X:f(X)\equiv_T B_j\}
\]
are invariant. If none contained a cone, every complement would
contain a cone. Their countable intersection, further intersected with
the original cone, would still contain a cone, contradicting that the
$E_j$ cover it. Thus some $E_j$ contains a cone, proving Part I for
this class of functions.

In particular, if there is a fixed oracle $B$ with $f(X)\le_T B$
throughout a cone, then $f$ is constant in degree on a smaller cone.
Only countably many reals are computable from $B$, since oracle
programs are indexed by integers. The word ``fixed'' is essential.
The condition $f(X)\le_T X$ allows the bound to vary with $X$ and
does not put all output degrees below one oracle. Replacing the latter
condition by the former would bypass a central part of the conjecture.

\paragraph{What cone decisions say about comparison.}
For two invariant functions $f,g$, the set
$\{X:f(X)\le_Tg(X)\}$ is invariant and so is the reversed comparison
set. Deciding both on cones and intersecting gives one uniform pattern:
equality of degrees, strict inequality in one direction, or
incomparability at every point of a cone. Thus determinacy reduces
comparison to a stable alternative but does not itself remove the
incomparability alternative. Part II requires its removal for functions
above the identity, as well as well-foundedness of the resulting order.

The same observation with $g(X)=X$ leaves possible eventual regression
or incomparability unless one proves the eventual-constant conclusion
of Part I. Stating that every invariant predicate is decided on a cone
does not determine which side of that predicate must be selected.
This is the first gap in a proof that invokes only the cone theorem.

\paragraph{Jump is well-defined and strictly raises the cone class.}
If $U\le_TV$, a $V$-oracle can simulate a $U$-oracle computation,
so the oracle halting problems satisfy $U'\le_TV'$. Consequently
pointwise jump preserves Turing invariance and cone equivalence, and
it is order-preserving for $\le_M$. For every $U$, one has $U\le_TU'$
and $U'\not\le_TU$. The first reduction uses a program which halts
exactly when the queried bit of $U$ is one. For the second, a purported
$U$-decision procedure for its own halting set yields a $U$-program
which, on its own index, halts precisely when the procedure says it
does not halt. This is the usual diagonal contradiction, valid
relative to every oracle.

It follows pointwise, and therefore on every cone, that
\[
 f(X)<_T f(X)',\qquad [f]<_M[f'].
\]
Here strictness means $f\le_M f'$ but $f'\not\le_M f$; the latter
is impossible on even a single input by the diagonal argument.
If $f$ is above the identity, so is $f'$. These facts establish a
strictly increasing operation, not that it is the immediate successor.
The missing assertion is that no invariant $g$ can satisfy
\[
 [f]<_M[g]<_M[f'].
\]
Pointwise strictness does not rule out such an intermediate function,
nor does it prove that arbitrary classes are comparable or that
descending chains are absent.

\paragraph{Tests of proposed exotic examples.}
For a fixed real $A$, the map $f_A(X)=X\oplus A$ looks different from
the identity but agrees with it in degree on the cone above $A$.
It produces no new Martin-order class. On the other hand, define
$h(X)=X$ when $0\notin X$ and $h(X)=X'$ when $0\in X$. This is
not Turing-invariant: toggling the first bit changes no Turing degree
of the input, but switches the output between that degree and its
strict jump. Therefore it is not a counterexample to the conjecture.
This test guards against constructing arbitrary case distinctions on
representatives rather than on degrees.

Finite iterates $X,X',X'',\ldots$ do give a strict chain of invariant
functions above the identity. Enumerating those examples does not
prove exhaustiveness, and assigning names to putative transfinite
iterates would not prove the specified well-order and successor law.

\paragraph{Literature, verification, and outcome.}
[S1] states the Turing-degree conjecture while studying an
enumeration-degree analogue; results about that analogue do not
automatically transfer. Lutz--Siskind [E1] establish Part I for the
additional order-preserving class and carefully distinguish the
choice assumptions of different proofs. The present arguments use
only the elementary cone proof and its countable completeness, with
the indicated consequence of determinacy. Verification is logical:
finite computations cannot establish oracle noncomputability or
quantification over all reals. The strongest partial conclusions are
eventual constancy for countable degree range and strictness of the
jump operation. General Part I, comparability and well-foundedness in
Part II, and the immediate-successor property remain
\textbf{UNRESOLVED}.

\subsection{Sources}
\begin{itemize}
\item[S1]Antonio Nakid Cordero, Martin's Conjecture in the Enumeration Degrees (2025), Introduction, Conjecture 1.2 and following paragraph.
\url{https://arxiv.org/html/2510.19147}.
\textit{Formulation source. Consulted 24 September 2026: Martin conjecture Parts I and II, cone order, and jump successors.}
\item[E1]Patrick Lutz and Benjamin Siskind, \emph{Part 1 of Martin's
Conjecture for order-preserving and measure-preserving functions},
arXiv:2305.19646, Sections 1.4 and 4.3.
\url{https://arxiv.org/html/2305.19646}.
\textit{Primary research source for the order-preserving subcase and
the distinction between countable choice for reals (a consequence of
AD) and stronger choice assumptions in other arguments. Consulted
27 September 2026; no stronger axiom is silently added here.}

\end{itemize}
\end{document}
